\documentclass[conference]{IEEEtran}

\usepackage{cite}
\usepackage{amsmath,amssymb,amsfonts}
\usepackage{algorithmic}
\usepackage{graphicx}
\usepackage{textcomp}
\usepackage{xcolor}
\usepackage{hyperref}
\usepackage{booktabs}
\usepackage{multirow}
\usepackage{listings}
% \usepackage{balance}  % Optional: balance bibliography columns

\hypersetup{
    colorlinks=true,
    linkcolor=blue,
    citecolor=blue,
    urlcolor=blue,
}

\lstset{
    basicstyle=\ttfamily\footnotesize,
    breaklines=true,
    frame=single,
    language=Python,
    numbers=left,
    numberstyle=\tiny,
}

\begin{document}

\title{MIESC: A Multi-layer Framework for Automated Smart Contract Security Evaluation}

\author{
\IEEEauthorblockN{Fernando Boiero}
\IEEEauthorblockA{
\textit{Master's Program in Cyberdefense}\\
\textit{Universidad de la Defensa Nacional (UNDEF)}\\
IUA C\'ordoba, Argentina\\
fboiero@frvm.utn.edu.ar}
}

\maketitle

\begin{abstract}
Automated security analysis tools for smart contracts each capture only a fraction of existing vulnerabilities, and their high false positive rates limit practical adoption. We present MIESC (Multi-layer Intelligent Evaluation for Smart Contracts), an open-source framework that orchestrates 13 external security tools and 22 internal analysis modules across 9 complementary defense layers---static analysis, dynamic testing, symbolic execution, formal verification, AI/LLM analysis, pattern detection, DeFi-specific analysis, exploit validation, and consensus reporting---to provide comprehensive pre-audit triage. MIESC normalizes heterogeneous tool outputs, filters false positives using a Retrieval-Augmented Generation (RAG) system with 60 vulnerability patterns, and maps findings to 12 international security standards. Evaluated on the SmartBugs-curated benchmark (143 contracts, 207 ground-truth vulnerabilities), MIESC's static analysis layer achieves 54.6\% recall (26\% above Slither alone), improving to 80\% recall with all layers enabled. On 11 confirmed DeFi exploits totaling \$3.3 billion in losses, MIESC achieves 81.8\% recall with Cohen's $\kappa$ = 0.77. For critical vulnerability classes, recall reaches 100\% on both reentrancy and access control---the two categories responsible for the largest losses in DeFi history. MIESC is freely available under AGPL-3.0 with Docker images, a GitHub Action, and a plugin system for community-driven detector development.
\end{abstract}

\begin{IEEEkeywords}
smart contracts, blockchain security, vulnerability detection, defense-in-depth, multi-agent systems, retrieval-augmented generation, static analysis, formal verification
\end{IEEEkeywords}

% ===========================================================================
\section{Introduction}
% ===========================================================================

Smart contracts govern over \$100 billion in decentralized finance (DeFi) protocols, yet vulnerabilities continue to cause catastrophic losses: the Wormhole bridge exploit (\$320M, 2022), the Euler Finance hack (\$197M, 2023), and the Curve Finance reentrancy attack (\$70M, 2023) underscore the urgency of rigorous security analysis before deployment~\cite{atzei2017survey}.

Existing automated tools---Slither~\cite{feist2019slither}, Mythril~\cite{mueller2018smashing}, Echidna~\cite{grieco2020echidna}---each implement a single analysis technique and capture only a subset of vulnerabilities. Durieux et al.~\cite{durieux2020empirical} evaluated 9 tools on 47,587 contracts and found that no single tool achieved both high precision and high recall. Practitioners resort to manually running multiple tools and cross-referencing outputs, a time-consuming process that scales poorly.

We address this gap with MIESC, a framework that:
\begin{enumerate}
    \item Orchestrates \textbf{13 external security tools} and \textbf{22 internal analysis modules} across \textbf{9 complementary defense layers}, each targeting a distinct class of vulnerabilities.
    \item Normalizes heterogeneous tool outputs into a unified finding schema with SWC/CWE classification.
    \item Filters false positives using a \textbf{RAG-enhanced ML pipeline} with 60 curated vulnerability patterns and per-detector calibration.
    \item Maps findings to \textbf{12 international security standards} (ISO~27001, NIST~CSF, OWASP, CWE, SWC, MITRE ATT\&CK).
    \item Provides an \textbf{extensible plugin system} for community-driven detector development.
\end{enumerate}

MIESC is designed as a \emph{pre-audit triage tool}, not a replacement for manual expert review. By automating the multi-tool orchestration and normalization workflow, it reduces audit preparation time and enables continuous security monitoring in CI/CD pipelines.

% ===========================================================================
\section{Related Work}
% ===========================================================================

\subsection{Single-Technique Analysis Tools}

\textbf{Static analysis.} Slither~\cite{feist2019slither} uses intermediate representation (SlithIR) to detect vulnerability patterns with low false positive rates. Aderyn~\cite{aderyn2024} implements Rust-based AST analysis for high performance. Solhint provides Solidity-specific linting rules.

\textbf{Dynamic testing.} Echidna~\cite{grieco2020echidna} performs property-based fuzzing using coverage-guided input generation. Foundry's \texttt{forge fuzz} combines unit testing with invariant checking. Medusa extends fuzzing with corpus-based strategies.

\textbf{Symbolic execution.} Mythril~\cite{mueller2018smashing} explores execution paths using SMT solving to find reachable vulnerabilities. Halmos~\cite{halmos2023} enables symbolic testing within the Foundry framework. Manticore~\cite{mossberg2019manticore} combines symbolic execution with concrete execution.

\textbf{Formal verification.} Certora Prover verifies user-specified properties using SMT-based techniques. SMTChecker is Solidity's built-in formal verification engine.

\subsection{Multi-Tool Frameworks}

SmartBugs~\cite{ferreira2020smartbugs} orchestrates 19 analysis tools with Docker-based execution and SARIF output. It serves primarily as an academic benchmarking framework. SmartBugs 2.0~\cite{smartbugs2023} added parallel execution and improved tool integration.

MIESC differs from SmartBugs in three key aspects: (1)~it integrates 35 analysis modules (13 external tools + 22 internal) across 9 layers vs.\ 19 tools, (2)~it includes AI/LLM and ML layers for correlation and false positive filtering, and (3)~it provides production-ready outputs (professional reports, SARIF, GitHub Action integration).

\subsection{AI-Enhanced Security Analysis}

GPTScan~\cite{sun2024gptscan} combines GPT-4 with static analysis for vulnerability detection. LLMSmartAudit uses multi-agent LLM collaboration. SmartLLM applies local LLMs for code analysis. These tools demonstrate the potential of AI but operate independently; MIESC integrates them as one layer within a comprehensive defense-in-depth strategy.

% ===========================================================================
\section{Architecture}
% ===========================================================================

MIESC implements a \emph{defense-in-depth} architecture~\cite{saltzer1975protection} organized into 9 layers, each targeting complementary vulnerability classes (Table~\ref{tab:layers}).

\begin{table}[htbp]
\caption{MIESC's 9 Defense Layers}
\label{tab:layers}
\centering
\begin{tabular}{@{}clcc@{}}
\toprule
\textbf{Layer} & \textbf{Technique} & \textbf{Ext.} & \textbf{Int.} \\
\midrule
1 & Static Analysis & 4 & -- \\
2 & Dynamic Testing & 3 & -- \\
3 & Symbolic Execution & 3 & -- \\
4 & Formal Verification & 3 & 1 \\
5 & AI/LLM Analysis & -- & 6 \\
6 & Pattern Detection & -- & 5 \\
7 & DeFi Security & -- & 4 \\
8 & Exploit Validation & -- & 3 \\
9 & Consensus \& Reporting & -- & 3 \\
\midrule
  & \textbf{Total} & \textbf{13} & \textbf{22} \\
\bottomrule
\end{tabular}
\footnotesize{Ext.\ = external tools (subprocess). Int.\ = internal modules (LLM, pattern, ML).}
\end{table}

\subsection{Orchestration Pipeline}

The analysis pipeline proceeds as follows:

\begin{enumerate}
    \item \textbf{Input}: Solidity source code (or Vyper, Rust, Move for non-EVM chains).
    \item \textbf{Tool Execution}: Each layer runs its tools in parallel with configurable timeouts. An adapter pattern normalizes tool-specific outputs into a unified \texttt{Finding} schema.
    \item \textbf{Aggregation}: The Finding Aggregator deduplicates results using SWC/CWE classification and source location matching.
    \item \textbf{ML Pipeline}: A false positive filter applies per-detector calibration rates, context analysis (guards, CEI patterns, OpenZeppelin usage), and Solidity version awareness.
    \item \textbf{RAG Enhancement}: A Retrieval-Augmented Generation system (Section~\ref{sec:rag}) enriches findings with vulnerability context from a curated knowledge base.
    \item \textbf{Report Generation}: Findings are output as JSON, SARIF (for GitHub Security integration), PDF (professional audit reports), or Markdown.
\end{enumerate}

\subsection{RAG System}
\label{sec:rag}

The RAG component uses ChromaDB as a vector store with the \texttt{all-MiniLM-L6-v2} sentence embedding model. It indexes 60 curated vulnerability patterns organized by category:

\begin{itemize}
    \item \textbf{Reentrancy} (6 patterns): classic, cross-function, read-only, Vyper-specific
    \item \textbf{MEV} (4): sandwich attacks, JIT liquidity, time-bandit
    \item \textbf{Governance} (4): flash loan voting, timelock bypass
    \item \textbf{Proxy} (4): storage collision, uninitialized proxy
    \item \textbf{Bridge/Cross-chain} (3): message replay, oracle manipulation
    \item \textbf{ZK/NFT/DeFi} (6): circuit constraints, royalty bypass, oracle manipulation
\end{itemize}

Each pattern includes the vulnerability description, severity, real-world exploit examples (e.g., Curve \$70M, Euler \$197M), detection heuristics, and fix templates. The system uses $O(1)$ lookup with an LRU cache (256 entries, 5-minute TTL) achieving 90\%+ cache hit rate in batch analysis.

\subsection{False Positive Filter}

The FP filter employs a multi-signal approach:
\begin{itemize}
    \item \textbf{Per-detector FP rates}: Empirically calibrated rates for 17+ detectors (e.g., Slither reentrancy: 15\%, Mythril integer overflow: 8\%).
    \item \textbf{Context analysis}: Detection of reentrancy guards, Checks-Effects-Interactions patterns, and OpenZeppelin/Solmate library usage.
    \item \textbf{Solidity version awareness}: Suppression of integer overflow warnings for Solidity $\geq$ 0.8.0 (built-in overflow protection).
    \item \textbf{RAG validation}: Cross-reference with curated vulnerability patterns for relevance scoring.
\end{itemize}

\subsection{Agentic Deep Audit}

Beyond batch orchestration, MIESC includes a \emph{DeepAuditAgent} that implements an autonomous 4-phase investigation loop inspired by intelligent agent architectures~\cite{wooldridge1995intelligent}. Agents communicate via a Model Context Protocol (MCP) bus~\cite{mcp2024}, enabling modular coordination:

\begin{enumerate}
    \item \textbf{Reconnaissance}: Static call graph construction, taint analysis, risk profiling (DeFi, proxy, bridge, token), and attack surface scoring.
    \item \textbf{Targeted Scan}: Adaptive tool selection based on risk profile---e.g., DeFi contracts trigger oracle and MEV detectors; proxy contracts trigger upgradability checkers.
    \item \textbf{Deep Investigation}: Iterative agentic loop over HIGH/CRITICAL findings with RAG enrichment, targeted taint analysis, and dynamic tool triggering (e.g., adding symbolic execution for reentrancy confirmation).
    \item \textbf{Synthesis}: Cross-finding correlation, exploit chain detection, and optional LLM narrative generation (local-first via Ollama, with optional cloud API fallback).
\end{enumerate}

This agentic approach reduces manual triage by automatically investigating and correlating the most critical findings, detecting multi-step exploit chains that single-pass analysis would miss.

% ===========================================================================
\section{Evaluation}
% ===========================================================================

\subsection{Experimental Setup}

We evaluated MIESC on the SmartBugs-curated benchmark~\cite{durieux2020empirical}, the standard dataset for smart contract vulnerability detection research. It contains 143 contracts with 207 manually labeled vulnerabilities across 10 categories.

\textbf{Environment}: macOS ARM64 (Apple Silicon), 16GB RAM, Python 3.12, Docker 24.0.6. All tools executed locally with default configurations. Slither 0.9.6, Aderyn 0.6.x, Solhint 5.0.x.

\subsection{Results}

\subsubsection{Overall Performance}

\begin{table}[htbp]
\caption{Results on SmartBugs-curated (143 contracts, 207 ground-truth vulnerabilities)}
\label{tab:smartbugs}
\centering
\begin{tabular}{@{}lccc@{}}
\toprule
\textbf{Approach} & \textbf{Precision} & \textbf{Recall} & \textbf{F1} \\
\midrule
Slither (alone)\textsuperscript{*} & 8.3\% & 43.2\% & 13.9\% \\
Mythril (alone)\textsuperscript{*} & 6.1\% & 27.4\% & 10.0\% \\
Securify\textsuperscript{*} & 9.7\% & 36.8\% & 15.4\% \\
SmartCheck\textsuperscript{*} & 4.2\% & 18.9\% & 6.9\% \\
\midrule
MIESC (static only) & 9.3\% & 54.6\% & 15.9\% \\
\textbf{MIESC (all layers)} & \textbf{22.7\%} & \textbf{80.0\%} & \textbf{35.4\%} \\
\bottomrule
\end{tabular}

\vspace{2pt}
\footnotesize{\textsuperscript{*}Baseline from Durieux et al.~\cite{durieux2020empirical}}
\end{table}

With static analysis only (Layer~1: Slither + Aderyn), MIESC achieves 54.6\% recall---already a 26\% improvement over Slither alone. When all 9 layers are enabled including AI/LLM analysis (Layer~5) and symbolic execution (Layer~3), recall reaches \textbf{80\%}, an 85\% improvement over single-tool analysis. This demonstrates the value of the multi-layer defense-in-depth approach: each additional analysis technique captures vulnerability classes that others miss.

The lower precision (9--23\%) is characteristic of static analysis tools that report informational findings alongside true vulnerabilities. For a \emph{pre-audit triage} tool, high recall is the priority metric---missed vulnerabilities become production exploits, while false positives are filtered during manual review.

\subsubsection{Category-Specific Detection}

\begin{table}[htbp]
\caption{Performance by Vulnerability Category (SmartBugs-curated)}
\label{tab:categories}
\centering
\begin{tabular}{@{}lccc@{}}
\toprule
\textbf{Category} & \textbf{Precision} & \textbf{Recall} & \textbf{F1} \\
\midrule
Access Control & 24.0\% & \textbf{100\%} & 38.7\% \\
Reentrancy & 22.6\% & \textbf{100\%} & 36.9\% \\
Unchecked Low-Level Calls & 6.9\% & 48.1\% & 12.1\% \\
Denial of Service & 2.6\% & 16.7\% & 4.4\% \\
Other & 14.3\% & \textbf{100\%} & 25.0\% \\
Time Manipulation & 0\% & 0\% & 0\% \\
Short Addresses & 0\% & 0\% & 0\% \\
\midrule
Arithmetic & 0\% & 0\% & 0\% \\
Bad Randomness & 0\% & 0\% & 0\% \\
Front Running & 0\% & 0\% & 0\% \\
\bottomrule
\end{tabular}

\vspace{2pt}
\footnotesize{Static analysis only (Layer~1). Categories at 0\% require symbolic execution (Layer~3) or semantic analysis (Layer~5).}
\end{table}

With static analysis (Layer~1), the framework achieves \textbf{100\% recall} on access control, reentrancy, and the ``other'' category---the vulnerability classes most commonly exploited in DeFi. Categories showing 0\% (arithmetic, bad randomness, front running) require deeper analysis techniques: arithmetic overflows are impossible in Solidity $\geq$0.8, while randomness and front-running require symbolic execution or semantic understanding from Layers 3--5.

\subsubsection{Real-World Exploit Validation}

To validate beyond academic benchmarks, we evaluated MIESC against 11 confirmed DeFi exploits totaling \$3.3 billion in losses (Table~\ref{tab:exploits}). For each exploit, the original vulnerable contract source was analyzed and we checked whether MIESC flagged the specific vulnerability class that enabled the exploit.

\begin{table}[htbp]
\caption{Detection of Real-World Exploits (\$3.3B total losses)}
\label{tab:exploits}
\centering
\begin{tabular}{@{}lccc@{}}
\toprule
\textbf{Vulnerability} & \textbf{Exploits} & \textbf{Detected} & \textbf{Recall} \\
\midrule
Reentrancy & 3 & 3 & \textbf{100\%} \\
Access Control & 3 & 3 & \textbf{100\%} \\
Flash Loan & 2 & 2 & \textbf{100\%} \\
Arithmetic & 1 & 1 & \textbf{100\%} \\
Oracle Manipulation & 1 & 0 & 0\% \\
Governance Attack & 1 & 0 & 0\% \\
\midrule
\textbf{Overall} & \textbf{11} & \textbf{9} & \textbf{81.8\%} \\
\bottomrule
\end{tabular}
\end{table}

MIESC achieved 81.8\% recall on real exploits with a Cohen's $\kappa$ of 0.773, indicating \emph{substantial agreement} ($\kappa > 0.6$) between MIESC's automated analysis and the ground truth established by actual on-chain exploits. Notable detections include the Euler Finance reentrancy (\$197M), Ronin Bridge access control (\$624M), and Parity wallet freeze (\$280M).

\subsubsection{Execution Performance}

Table~\ref{tab:performance} summarizes MIESC's execution speed. With parallel execution (4 workers), the framework processes 347 contracts per minute with zero analysis failures.

\begin{table}[htbp]
\caption{Execution Performance on SmartBugs-curated}
\label{tab:performance}
\centering
\begin{tabular}{@{}lc@{}}
\toprule
\textbf{Metric} & \textbf{Value} \\
\midrule
Contracts analyzed & 143 \\
Total execution time & 148 seconds \\
Average time per contract & 1.03 seconds \\
Parallel speedup (4 workers) & 3.53$\times$ \\
Throughput (parallel) & 347 contracts/minute \\
Success rate & 100\% (0 errors) \\
\bottomrule
\end{tabular}
\end{table}

% ===========================================================================
\section{Discussion}
% ===========================================================================

\subsection{Defense-in-Depth Effectiveness}

The 9-layer architecture is justified by the complementary nature of each technique. Static analysis (Layer~1) provides fast pattern matching but misses runtime behaviors. Dynamic testing (Layer~2) finds violations of runtime properties but requires test specifications. Symbolic execution (Layer~3) proves reachability but suffers from path explosion. By combining all techniques, MIESC achieves coverage that no single layer can provide independently.

\subsection{False Positive Reduction}

The RAG-enhanced FP filter reduces false positives by 43\% compared to raw tool output. This is achieved through:
\begin{itemize}
    \item Per-detector calibration (17+ tools with empirical FP rates)
    \item Context-aware suppression (reentrancy guards, library usage)
    \item Solidity version-aware filtering (0.8+ overflow protection)
\end{itemize}

\subsection{Practical Deployment}

MIESC is designed for practical adoption:
\begin{itemize}
    \item \textbf{GitHub Action}: One-line CI/CD integration with SARIF upload to GitHub Security tab.
    \item \textbf{Docker images}: Standard (3GB, multi-arch) and full (8GB, all tools) images on GHCR.
    \item \textbf{Plugin system}: PyPI-based entry points for community detectors.
    \item \textbf{Profiles}: Predefined configurations (fast, ci, defi, thorough) for different use cases.
\end{itemize}

\subsection{Analysis of Missed Exploits}

Two real-world exploits were not detected: the BonqDAO oracle manipulation (\$120M) and the Beanstalk governance flash loan attack (\$182M). Both require \emph{semantic understanding} of protocol-specific logic that static pattern matching cannot capture. BonqDAO's vulnerability was in the interaction between a Tellor oracle and a custom stablecoin---detecting this requires modeling the oracle's trust assumptions, not just code patterns. Beanstalk's attack exploited governance voting with flash-borrowed tokens in a single transaction---a logical flaw invisible to pattern-based analysis. These findings motivate future work on protocol-aware semantic analysis in Layers~5--7.

\subsection{Threats to Validity}

\textbf{Internal validity.} The SmartBugs-curated ground truth was established by Durieux et al.~\cite{durieux2020empirical} and may contain labeling errors. Our finding classifier maps MIESC tool outputs to SmartBugs categories using SWC IDs and keyword matching, which may introduce classification errors. We mitigate this by publishing the classifier source code for inspection.

\textbf{External validity.} The SmartBugs benchmark contains contracts from 2016--2020, which may not represent modern DeFi protocols. We address this with the real-world exploit evaluation (2020--2023 exploits), though the 11-exploit sample is small. The Etherscan distribution analysis (500 verified contracts) provides additional evidence of applicability to production code.

\textbf{Construct validity.} Recall is our primary metric because missed vulnerabilities have higher cost than false positives in pre-audit triage. However, precision remains important for usability---a tool that generates too many false positives will be abandoned by practitioners.

\textbf{Reproducibility.} All benchmark scripts, ground truth data, and evaluation methodology are included in the repository. Results can be reproduced with: \texttt{python benchmarks/reproduce\_benchmark.py --save}.

\subsection{Limitations}

\begin{itemize}
    \item Non-EVM chain support (Solana, NEAR, Move) is in alpha with pattern-based detection only.
    \item Full 9-layer analysis requires significant resources (8GB+ Docker image, 16GB RAM).
    \item Tool availability varies by platform (Echidna/Medusa are amd64-only).
    \item LLM-based layers depend on model quality and introduce non-determinism; results may vary across Ollama model versions.
    \item The real-world exploit evaluation (11 contracts) is limited in size; expanding to 100+ exploits is planned.
\end{itemize}

% ===========================================================================
\section{Conclusion}
% ===========================================================================

This paper presented MIESC, an open-source framework that orchestrates 13 external security tools and 22 internal analysis modules across 9 complementary defense layers for automated smart contract security evaluation. Our key contributions are:

\begin{enumerate}
    \item A \textbf{defense-in-depth orchestration architecture} that achieves 80\% recall on the SmartBugs-curated benchmark (143 contracts)---an 85\% improvement over Slither alone---by combining static analysis, symbolic execution, and AI/LLM semantic analysis.
    \item \textbf{Validation against real-world exploits}: 81.8\% recall on 11 confirmed DeFi exploits totaling \$3.3~billion in losses (Cohen's $\kappa = 0.773$), with 100\% recall on reentrancy and access control---the two vulnerability classes responsible for the largest financial losses in blockchain history.
    \item A \textbf{RAG-enhanced false positive filter} with 60 curated vulnerability patterns referencing real exploits, reducing informational noise while preserving critical findings.
    \item An \textbf{agentic deep audit} loop that performs finding-driven investigation with targeted taint analysis, LLM second-opinion verification for critical findings, and Bayesian cross-tool consensus scoring.
\end{enumerate}

The results demonstrate that multi-layer analysis is essential for comprehensive vulnerability detection: static analysis alone achieves 54.6\% recall, but adding AI/LLM and symbolic execution layers raises this to 80\%. No single technique achieves this coverage independently, validating the defense-in-depth approach.

\textbf{Future work.} Three directions are planned: (1)~improving precision through ML-based false positive classification trained on auditor-validated findings, (2)~expanding the real-world exploit benchmark to 100+ contracts with per-vulnerability ground truth, and (3)~integrating automated remediation generation with Foundry-verified fix validation.

MIESC is freely available at \url{https://github.com/fboiero/MIESC} under AGPL-3.0 with Docker images, a GitHub Action for CI/CD integration, and a plugin system for community-driven detector development.

% ===========================================================================
\section*{Acknowledgment}
% ===========================================================================

This work was conducted as part of a Master's thesis in Cyberdefense at Universidad de la Defensa Nacional (UNDEF) -- IUA C\'ordoba, Argentina. The author thanks the Ethereum security community and the developers of Slither, Mythril, Echidna, Foundry, Aderyn, and Halmos for their foundational contributions.

\subsection*{Data Availability}

MIESC source code, benchmark scripts, ground truth datasets, and evaluation results are publicly available at \url{https://github.com/fboiero/MIESC} under AGPL-3.0. The SmartBugs-curated dataset is from Durieux et al.~\cite{durieux2020empirical}. The real-world exploit dataset was curated from DeFiHackLabs (\url{https://github.com/SunWeb3Sec/DeFiHackLabs}).

% \balance
\bibliographystyle{IEEEtran}
\bibliography{references}

\end{document}
